\section{全核波形中的点积事务}
为观察点积在全核中的输入、响应和提交，2026年11月5日使用稳定八项点积版本的clean-02封存生成RTL，单独构建了支持FST的EMU。1681项生成输入与原清单一致；该模型运行原functional-2镜像，加载正常NEMU参考库并启用DiffTest，最终到达GOOD TRAP、退出码0。以下波形将同一笔正数点积从执行输入关联到架构提交。

\subsection{源数据、结果与身份}
两个源寄存器均为\seqsplit{0xa5a5a5a5a5a5a5a50807060504030201}，低八个有符号字节依次为1至8。预期结果为
\[
1^2+2^2+\cdots+8^2=204.
\]
实际128位响应为\seqsplit{0x000000000000000000000000000000cc}，符合低32位保存结果、高96位清零的布局。高64位A5哨兵同时检验了源元素选择范围。

{\small\begin{longtable}{p{.24\linewidth}p{.68\linewidth}}\toprule
关联字段 & 观测值\\\midrule
ROB与物理目的 & ROB0:39；物理目的103（十六进制67）\\
提交观察 & \texttt{io\_valid=1}、\texttt{io\_vecwen=1}；架构目的v10\\
PC与机器字 & PC=\texttt{0x80001b8c}；机器字\texttt{0x0294050b}\\
\bottomrule\end{longtable}}

\subsection{从有效区间到握手边沿}
波形中的valid表示一段有效区间，事务接受则发生在valid与ready同时有效的时钟上升沿。两者按下表区分，便于与日志中的周期坐标对照。

{\small\begin{longtable}{p{.25\linewidth}p{.17\linewidth}p{.48\linewidth}}\toprule
事件 & 显示坐标（ps） & 观察依据\\\midrule
输入有效（A） & 678694 & 输入valid、ready均为1，源值及身份有效\\
输入接受边沿 & 678696 & 边沿前握手有效，边沿后P槽有效\\
输出有效（B） & 678698 & R槽有效，输出valid、ready均为1，结果204\\
输出握手边沿 & 678700 & 边沿前输出有效，边沿后R槽清空\\
DiffTest提交观察（C） & 678710 & 提交有效，PC、编码、ROB与目的相符\\
\bottomrule\end{longtable}}

同一事务的O2日志记录接受、响应和内部ROB提交周期分别为339296、339298、339301。接受到响应相隔两个全局模型周期，与第二阶段的两槽时序一致。内部ROB提交对应模型坐标678706，DiffTest观察口在678710有效，反映两拍观察流水延迟。分析执行时延时使用实际握手和内部提交坐标，分析DiffTest接口时使用其观察坐标。

图上方为全局时钟，下方为点积局部门控时钟；局部空闲时可停止时钟。显示单位ps是仿真模型的时间坐标，相邻全局周期相隔两个坐标单位。程序性能采用DUT周期，物理时序由综合和STA评估。

以下三幅GTKWave截图分别展示输入、响应和提交。截图保留信号值，仅裁去层级栏和外部空白；完整原图、事务采样、运行回执与文件摘要见\seqsplit{evidence/vdot-wave-20261105/}和\seqsplit{figures/source-screenshots/}。

\begin{landscape}
\subsection{输入与接受}
\reportfigure{S3-WAVE-01}{正数点积的输入与接受边沿}{10.4cm}{游标678694处输入valid、ready为1；源低八字节为1至8，ROB值27和物理目的67按十六进制显示。678696上升沿接收请求，P槽有效。}{2026年11月5日八项点积全核trace原始截图。}
\clearpage
\subsection{结果与响应}
\reportfigure{S3-WAVE-02}{完整128位结果与响应握手}{10.4cm}{游标678698处R槽及输出valid、ready有效，结果为204且高96位为零；678700边沿完成握手并释放R槽。此时提交观察口io\_valid为0。}{2026年11月5日八项点积全核trace原始截图。}
\clearpage
\subsection{提交与架构目的}
\reportfigure{S3-WAVE-03}{DiffTest提交观察与架构目的}{10.4cm}{游标678710处提交有效，PC80001B8C、机器字0294050B、ROB027、物理目的67、架构目的0A（v10）及向量写使能对应同一事务。观察口的两拍流水延迟已在正文对齐。}{2026年11月5日八项点积全核trace原始截图。}
\end{landscape}
